\begin{figure}[ht]
\centering
\begin{tikzpicture}[x=1.1cm, y=0.85cm, font=\small]

% Axes
\draw[->, thick] (0, 0) -- (10.6, 0)
  node[anchor=west, font=\scriptsize] {iteration $k$};
\draw[->, thick] (0, 0) -- (0, 6.6)
  node[anchor=south, font=\scriptsize] {$\log_{10}\norm{\vect{r}_k}_2$};

% x ticks
\foreach \k in {0,2,4,6,8,10} {
  \draw (\k, 0.08) -- (\k, -0.08) node[anchor=north, font=\scriptsize] {\k};
}
% y ticks: 0 down to -6, drawn upward as 0..6 then label negated
\foreach \y/\lab in {0/0, 1/-1, 2/-2, 3/-3, 4/-4, 5/-5, 6/-6} {
  \draw (0.08, \y) -- (-0.08, \y) node[anchor=east, font=\scriptsize] {\lab};
}

% Jacobi: slow linear decay (spectral radius ~ 0.8 => ~0.1 decade/iter)
\draw[very thick, engamber]
  (0, 0) -- (1, 0.32) -- (2, 0.64) -- (3, 0.97) -- (4, 1.29)
  -- (5, 1.61) -- (6, 1.93) -- (7, 2.25) -- (8, 2.58) -- (9, 2.90) -- (10, 3.22);
\node[engamber, anchor=west] at (7.1, 2.6) {Jacobi};

% Gauss-Seidel: roughly twice the rate of Jacobi here
\draw[very thick, engblue]
  (0, 0) -- (1, 0.62) -- (2, 1.25) -- (3, 1.87) -- (4, 2.50)
  -- (5, 3.12) -- (6, 3.74) -- (7, 4.37) -- (8, 4.99) -- (9, 5.62) -- (10, 6.24);
\node[engblue, anchor=south] at (8.2, 5.3) {Gauss--Seidel};

% Conjugate gradient: superlinear, faster late
\draw[very thick, engred]
  (0, 0) -- (1, 0.55) -- (2, 1.30) -- (3, 2.25) -- (4, 3.30)
  -- (5, 4.40) -- (6, 5.55) -- (7, 6.40);
\node[engred, anchor=west] at (4.4, 4.0) {conjugate gradient};

\node[anchor=north, font=\scriptsize, text width=10.5cm, align=center]
  at (5.0, -1.0) {
  Residual norm against iteration for a symmetric positive-definite
  system, schematic. Stationary iterations (Jacobi, Gauss--Seidel)
  drop the residual by a fixed factor per step set by the iteration's
  spectral radius. Conjugate gradient accelerates as it builds the
  Krylov subspace and terminates in at most $n$ steps in exact
  arithmetic.
};

\end{tikzpicture}
\caption[Convergence of stationary and Krylov iterations]{Schematic
residual-norm history for three iterative solvers on a symmetric
positive-definite system. The vertical axis is $\log_{10}$ of the
residual two-norm, so a straight descending line is geometric
convergence at a fixed rate per iteration. Jacobi and Gauss--Seidel
each converge linearly, with Gauss--Seidel typically about twice as
fast on the same matrix; conjugate gradient converges superlinearly,
its rate governed by $\sqrt{\kappa(\mat{A})}$ rather than
$\kappa(\mat{A})$, which is the reason a good preconditioner that
lowers $\kappa$ pays off so steeply. The slopes are illustrative; the
exact rates depend on the matrix spectrum.}
\label{fig:vol02:ch09:iterative-convergence}
\end{figure}
